\documentclass[10pt,a4paper]{amsart}
\usepackage[margin=19mm]{geometry}
\usepackage{amsmath,amssymb,mathtools}
\usepackage[scr=rsfs,cal=euler]{mathalfa}
\usepackage{xfrac}
\usepackage[unicode,hidelinks,bookmarks=false]{hyperref}
\newcommand{\ie}{i.e.\ }
\newcommand{\cf}{cf.\ }
\newcommand{\resp}{respectively\ }
\newcommand{\Subsection}{\subsection}
\providecommand{\nobreakdash}{\nobreak\hbox{-}}
\setlength{\emergencystretch}{2em}
\multlinegap0pt
\usepackage{lmodern}
\usepackage{mathtools}
\usepackage{xurl}
\newcommand{\Z}{\mathbb Z}
\newcommand{\F}{\mathbb F}
\DeclareMathOperator{\GL}{GL}
\DeclareMathOperator{\St}{St}
\DeclareMathOperator{\id}{id}
\DeclareMathOperator{\im}{im}
\DeclareMathOperator{\coker}{coker}
\DeclareMathOperator{\Hom}{Hom}
\DeclareMathOperator{\Ext}{Ext}
\DeclareMathOperator{\Tor}{Tor}
\newtheorem{theorem}{Theorem}[section]
\newtheorem{lemma}[theorem]{Lemma}
\newtheorem{proposition}[theorem]{Proposition}
\newtheorem{corollary}[theorem]{Corollary}
\newtheorem{definition}[theorem]{Definition}
\newtheorem{remark}[theorem]{Remark}
\numberwithin{equation}{section}
\newcommand{\eqnum}[1][]{%
  \refstepcounter{equation}%
  \if\relax\detokenize{#1}\relax\else\label{#1}\fi
  \leqno\hbox{\normalfont(\theequation)}}
\newcounter{displaytag}
\newcommand{\eqrownum}[1][]{%
  \hbox{\refstepcounter{equation}%
    \if\relax\detokenize{#1}\relax\else\label{#1}\fi
    \normalfont(\theequation)}}
\newcommand{\eqrowtag}[2]{%
  \hbox{\refstepcounter{displaytag}%
    \def\@currentlabel{#1}%
    \label{#2}%
    \normalfont(#1)}}
\hypersetup{pdfauthor={Huynh Viet Khanh},pdftitle={General linear and Steinberg groups over the Leavitt algebra L\_F2(1,2)}}
\begin{document}
\title[General linear and Steinberg groups over the Leavitt algebra]{General linear and Steinberg groups over the Leavitt algebra $L_{\mathbb F_2}(1,2)$}
\author{280 An Duong Vuong Str., Ho Chi Minh City, Vietnam; Huynh Viet Khanh}
\date{}
\maketitle
\noindent\emph{Source and licence.} General linear and Steinberg groups over the Leavitt algebra $L_{\mathbb F_2}(1,2)$. Version arXiv:2609.08428v2 (CC BY 4.0). This material is distributed under the Creative Commons Attribution 4.0 International licence, http://creativecommons.org/licenses/by/4.0/, as recorded in the arXiv OAI licence metadata for this exact version and shown on the abstract page https://arxiv.org/abs/2609.08428v2. Full rights notice: licenses/arxiv-2609.08428-CC-BY-4.0.txt.\par\smallskip
\noindent\textit{Editorial scope.} Counted unit: the original \texttt{theorem} environment \texttt{thm:raw-criterion} at source lines 453--462 together with its complete proof at lines 464--553; the delivered document keeps source lines 115--554 so that every referenced label and every citation resolves inside it. Native editable source is the paper's own concatenated TeX; the AMS wrapper is editorial formatting only. No publisher class, upstream script or unrelated artwork is redistributed.
\begin{lemma}\label{lem:leaf-coordinates}
For a complete ordered leaf set $(\alpha_1,\ldots,\alpha_m)$, let $L=(\alpha_1,\ldots,\alpha_m)$ and $L^*=(\alpha_1^*,\ldots,\alpha_m^*)^{\mathrm t}$. Then
$$
LL^*=1\qquad\text{and}\qquad L^*L=I_m. \eqnum[eq:leaf-identities]
$$
Consequently $x\mapsto Lx$ and $a\mapsto L^*a$ are mutually inverse maps of right modules $R^m\rightleftarrows R$, and
$$
M_m(R)\rightleftarrows R,\qquad Z\longmapsto LZL^*\qquad\text{and}\qquad a\longmapsto L^*aL \eqnum[eq:leaf-ring-isomorphisms]
$$
are mutually inverse unital ring isomorphisms. Every positive integer $m$ occurs as the size of a complete leaf set.
\end{lemma}

\begin{proof}
Distinct leaves are prefix-incomparable. Cancellation therefore implies $\alpha_i^*\alpha_j=\delta_{ij}$. Starting with $1=1\cdot1^*$, each leaf replacement preserves the sum of the range projections, because $\alpha\alpha^*=\alpha ee^*\alpha^*+\alpha ff^*\alpha^*$. Hence $\sum_i\alpha_i\alpha_i^*=1$, proving~\eqref{eq:leaf-identities}. These identities verify the two pairs of inverse maps and show that the ring maps preserve products and units. A complete leaf set of size $m$ is the comb
$$
L_m=(e^{m-1},e^{m-2}f,\ldots,ef,f),\qquad\text{with }L_1=(1).
$$
Its recurrence is $L_{m+1}=(eL_m,f)$.
\end{proof}

\begin{remark}\label{rem:proper-standard-V}
The standard copy of Thompson's group $V$ consists of the units $\sum_i\alpha_i\beta_i^*$ associated with bijections between complete ordered leaf sets of the same size. These units are unitary. In contrast, $u=1+ef^*$ satisfies $u^{-1}=u$ and $u^*=1+fe^*\ne u$, since $e^*(ef^*-fe^*)f=1$. Thus the standard copy of $V$ is a proper subgroup of $G$.
\end{remark}

The following proposition combines the stable calculation of \cite{ABC09} with an identity involving finite sums, as in the theory of sum rings \cite[\S2]{Wagoner72}.

\begin{proposition}\label{prop:zero-padding}
For every $r\geq1$, the standard inclusion $G_r\to G_{r+1}$, $g\mapsto\operatorname{diag}(g,1)$, induces zero on $H_n(-,\Z)$ for every $n>0$.
\end{proposition}

\begin{proof}
For the graph with one vertex and two loops, the algebraic $K$-theory sequence of Ara--Brustenga--Corti\~nas \cite[Theorem~7.6]{ABC09} has map $1-2=-1$ on $K(\F_2)$. The graph is row-finite and has no sinks. The coefficient field is regular supercoherent: every polynomial ring in finitely many variables over $\F_2$ is regular Noetherian. By the theorem, $K(R)$ is the homotopy cofiber of an equivalence. Hence $K_i(R)=0$ for all $i>0$.

Let $\GL_\infty(R)$ be the direct limit under standard inclusions. The plus construction $B\GL_\infty(R)^+$ has fundamental group $K_1(R)$ and higher homotopy groups $K_i(R)$, $i\geq2$, and preserves integral homology \cite[Chapter~IV, Definitions~1.1 and~1.1.1]{Weibel13}. It is a connected CW complex with trivial homotopy groups, hence is contractible by Whitehead's theorem. Thus
$$
H_n(\GL_\infty(R),\Z)=0\qquad\text{for }n>0. \eqnum[eq:stable-homology-zero]
$$

Define $c:G\to G$ by $c(u)=eue^*+ff^*$. Orthogonality implies $c(u)c(v)=c(uv)$, $c(1)=1$, and $c(u)^{-1}=c(u^{-1})$. For $s\geq0$, let $\Phi_s:R\to M_{2^s}(R)$ have entries $\Phi_s(a)_{\alpha,\beta}=\alpha^*a\beta$, where $\alpha,\beta$ range over the words of length $s$. These are the ring isomorphisms of Lemma~\ref{lem:leaf-coordinates}. Ordering the words at depth $s+1$ by $e\alpha$ followed by $f\alpha$, we have
$$
\Phi_{s+1}(c(u))=\operatorname{diag}(\Phi_s(u),I_{2^s}). \eqnum[eq:dyadic-compression]
$$

Thus the system $(G,c)$ is isomorphic to the dyadic cofinal subsystem of the standard general linear system. The normalized bar complex commutes with this filtered direct limit, since every generator and every finite chain occurs at a finite stage. According to~\eqref{eq:stable-homology-zero}, $\varinjlim(H_n(G,\Z),c_*)=0$. Therefore every element of $H_n(G,\Z)$ is killed by a positive iterate of $c_*$; that is, $c_*$ is locally nilpotent.

It remains to prove that $c_*$ is idempotent. The units associated with the complete leaf sets $(e,fe,ff)$ and $(ee,ef,f)$ are
$$
w=ee\,e^*+ef(fe)^*+f(ff)^*\qquad\text{and}\qquad w^{-1}=e(ee)^*+fe(ef)^*+ff\,f^*.
$$
Both inverse products follow from~\eqref{eq:leaf-identities}. Moreover $we=ee$ and $w(ff^*)w^{-1}=ef(ef)^*+ff^*$. Hence
$$
wc(u)w^{-1}=ee\,u(ee)^*+ef(ef)^*+ff^*=c(c(u)). \eqnum[eq:compression-idempotence]
$$
Inner conjugation acts trivially on homology, so $c_*^2=c_*$. Since $c_*$ is also locally nilpotent, it follows that $c_*=0$ on every positive homology group.

For $Z\in G_r$, the comb recurrence takes the form
$$
L_{r+1}\operatorname{diag}(Z,1)L_{r+1}^* =c(L_rZL_r^*). \eqnum[eq:padding-compression]
$$
Under the leaf-coordinate group isomorphisms, this identity identifies standard stabilization with $c$. Since $c_*=0$, the standard inclusion induces zero on positive homology.
\end{proof}

The map $c$ used here is a group homomorphism. Its affine formula sends $0$ to $ff^*\neq0$, so it is not a ring homomorphism. Proposition~\ref{prop:zero-padding} concerns the specified standard inclusions; the coordinate isomorphisms of Lemma~\ref{lem:leaf-coordinates} induce homology isomorphisms.

\section{Ordered frames and simultaneous extensions}
\label{sec:common-cones}

We use completable ordered frames over nonzero unital $\F_2$-algebras; compare the special unimodular frames of \cite[\S2]{NS90} and \cite[\S2]{Suslin84}. Since a free module may have bases of different finite cardinalities, we specify the number of basis vectors in each complement.

\begin{definition}\label{def:ordered-frames}
Let $A$ be such an algebra and write $\Gamma_r=\GL_r(A)$. The ordered frame semisimplicial set $X_r(A)$ has as its $p$-simplices the tuples $(v_1,\ldots,v_k)$, where $k=p+1\leq r$, admitting a decomposition
$$
A^r=v_1A\oplus\cdots\oplus v_kA\oplus C,\qquad\text{with }C\cong A^{r-k}. \eqnum[eq:frame-decomposition]
$$
Each map $A\to v_iA$, $a\mapsto v_i a$, must be an isomorphism. Since $A\neq0$, the condition $C\cong A^{r-k}$ implies that $C\neq0$ when $k<r$ and $C=0$ when $k=r$. The simplex records the ordered vectors and the formal integer $r-k$; the complement and its basis are not recorded. Faces delete vectors, and there are no simplices in dimensions $p\geq r$.
\end{definition}

Deleting a vector adds its summand to the complement, so the faces in Definition~\ref{def:ordered-frames} are well defined. The condition $C\cong A^{r-k}$ specifies a basis cardinality, which need not be an invariant of the module $C$.

To prove the homology vanishing required in Theorem~\ref{thm:acyclicity-criterion}, we extend each finite family of frames by a common transverse vector; compare \cite[\S2.7, Remark~4]{vdK80} and \cite[\S2, Lemmas~2.1--2.2]{Suslin84}. The Reduction Theorem for Leavitt path algebras treats one nonzero element \cite[Lemma~2.2.8, Theorem~2.2.11 and Corollary~2.2.12]{AAS17}. For $R$, we prove below a simultaneous version: one right multiplier works for a finite family, and the resulting left inverses have nonzero kernels with explicit finite bases.

If $\mu_1,\ldots,\mu_s$ are distinct positive words of lengths at most $D$ and $M>D$, then the words $\mu_i e^M f$ are pairwise prefix-incomparable. Indeed, for comparable original words write $\mu_j=\mu_i\delta$ with $\delta$ nonempty. If $\delta$ contains $f$, its first $f$ occurs before position $M+1$; if $\delta=e^a$, the first $f$ in $\delta e^M f$ occurs at position $M+a+1$. In either case $e^M f$ and $\delta e^M f$ disagree before the shorter word ends. Original words that were prefix-incomparable remain so. Therefore
$$
(\mu_i e^M f)^*(\mu_j e^M f)=\delta_{ij}. \eqnum[eq:word-separation]
$$
In particular, a nonempty sum of distinct positive words is nonzero in $R$: append $e^M f$ and multiply on the left by the star of one of the resulting words to obtain $1$.

\begin{lemma}\label{lem:word-multiplier}
Given finitely many nonzero $a_1,\ldots,a_s\in R$, there are $x\in R$ and nonempty positive words $\eta_1,\ldots,\eta_s$ such that $\eta_i^*a_ix=1$ for every $i$. Each decomposition $R=a_ixR\oplus\ker\eta_i^*$ has an explicit nonzero complement with a finite basis consisting of positive words.
\end{lemma}

\begin{proof}
Every element of $R$ is a finite sum of monomials $\alpha\beta^*$. To obtain this form, cancel every occurrence of a starred letter followed by an unstarred letter. Choose expressions of this form for the $a_i$, and choose $N$ at least as large as all lengths $|\beta|$ occurring in these expressions. Let $\gamma_1,\ldots,\gamma_t$ be all words of length $N$. Each $a_i\gamma_j$ is then a polynomial in positive words. For every $i$, at least one is nonzero, since $a_i=\sum_j(a_i\gamma_j)\gamma_j^*$, although the useful index $j$ may depend on $i$. Collect identical words over $\F_2$ in all these polynomials, and let $L$ bound the lengths of the remaining words.

Put $t_j=e^{j(L+1)}f$ and $x_0=\sum_j\gamma_jt_j$. The words contributed by $(a_i\gamma_j)t_j$ have lengths in $[j(L+1)+1,j(L+1)+1+L]$. These intervals are disjoint for distinct $j$, and concatenation by a fixed $t_j$ preserves distinct words within each interval. Thus $P_i=a_ix_0$ is a nonzero finite sum of distinct words in $e,f$ for every $i$.

We next make the support of each $P_i$ prefix-incomparable. Choose $M$ greater than every word length occurring in the $P_i$ and set $x=x_0e^Mf$. Equation~\eqref{eq:word-separation} makes the support of each $P_ie^Mf$ a prefix antichain. Choose $\eta_i$ in this support. Its coefficient is $1$, so $\eta_i^*a_ix=1$ by~\eqref{eq:word-separation}. Every $\eta_i$ is nonempty.

For a word $\eta=c_1\cdots c_h$, $h\geq1$, let $s_\ell=c_1\cdots c_{\ell-1}c'_\ell$, where $c'_\ell$ is the other binary letter. The words $\eta,s_1,\ldots,s_h$ form a complete leaf set. Hence
$$
\ker\eta^*=\bigoplus_{\ell=1}^h s_\ell R, \eqnum[eq:sibling-kernel]
$$
with inverse coordinate maps $(r_\ell)_\ell\mapsto\sum_\ell s_\ell r_\ell$ and $z\mapsto(s_\ell^*z)_\ell$. If $y=a_ix$ and $b=\eta_i^*$, then $by=1$, and
$$
R\oplus\ker b\rightleftarrows R,\qquad (t,z)\longmapsto yt+z\qquad\text{and}\qquad u\longmapsto(bu,u-ybu)
$$
are mutually inverse. The kernel in~\eqref{eq:sibling-kernel} is nonzero because $h\geq1$ and $s_\ell^*s_\ell=1$.
\end{proof}

\begin{proposition}\label{prop:frame-cone}
For a finite family of simplices of $X_r(R)$, each having at most $r-2$ frame vectors, there is a single vector $v$ extending every frame in the family to a simplex of $X_r(R)$. Each extended frame has a nonzero complement with an explicit basis of the cardinality required in Definition~\ref{def:ordered-frames}.
\end{proposition}

\begin{proof}
For an empty family, choose any standard vertex. Otherwise, let $W_i$ be the span of the $i$th frame, with its given ordered basis, let $k_i$ be its length, and put $q_i=r-k_i\geq2$.

Choose a completion $R^r=W_i\oplus C_i^0$ with $C_i^0\cong R^{q_i}$. Compose the quotient coordinate map with $L_{q_i}:R^{q_i}\to R$ to obtain a surjection $\rho_i:R^r\to R$ with kernel $W_i$. Let $\sigma_i:R\to R^r$ be the corresponding section with image $C_i^0$, so that $\rho_i\sigma_i=1$.

Fix $\Phi=L_r^*:R\to R^r$. The endomorphism $\rho_i\Phi$ of the right module $R$ is left multiplication by an element $a_i$. Since $\rho_i\Phi$ is surjective and $R\ne0$, one has $a_i\ne0$. Apply Lemma~\ref{lem:word-multiplier} to the $a_i$, and put $v=\Phi x$, $y_i=\rho_i v=a_ix$, and $b_i=\eta_i^*$. Since $b_iy_i=1$, the map $R\to R^r$ defined by $a\mapsto va$ is a split injection, with left inverse $b_i\rho_i$ for every $i$.

Put $w_i=v-\sigma_i y_i\in W_i$. We have
$$
R^r=W_i\oplus vR\oplus\sigma_i\ker b_i, \eqnum[eq:cone-decomposition]
$$
with coordinate map and inverse
$$
\begin{aligned}
 (w,t,z)&\longmapsto w+vt+\sigma_i z\qquad\text{and}\\
 u&\longmapsto\begin{pmatrix}
 (I-\sigma_i\rho_i)u-w_i b_i\rho_i u\\
 b_i\rho_i u\\
 \rho_i u-y_i b_i\rho_i u
 \end{pmatrix}.
 \end{aligned} \eqnum[eq:cone-coordinate-inverses]
$$
The first inverse component lies in $W_i$ because $\rho_i\sigma_i=1$ and $\rho_iw_i=0$; the last lies in $\ker b_i$ because $b_iy_i=1$. If $u=w+vt+\sigma_i z$, then $\rho_i u=y_it+z$ and $b_i\rho_i u=t$, so the inverse recovers $(w,t,z)$. Conversely, substituting the displayed components into the forward map returns $u+(v-w_i-\sigma_i y_i)b_i\rho_i u=u$.

Let $h_i=|\eta_i|$, and let $S_i=(s_1,\ldots,s_{h_i})$ be the row in~\eqref{eq:sibling-kernel}. The map
$$
\sigma_iS_iL_{h_i}^*L_{q_i-1}:R^{q_i-1} \longrightarrow\sigma_i\ker b_i \eqnum[eq:complement-formal-coordinates]
$$
is an isomorphism, with inverse on this complement $L_{q_i-1}^*L_{h_i}S_i^*\rho_i$. By $\rho_i\sigma_i=1$, $S_i^*S_i=I_{h_i}$, and~\eqref{eq:leaf-identities}, the inverse product on $R^{q_i-1}$ is the identity. The other product is $\sigma_iS_iS_i^*\rho_i$, which is the identity on $\sigma_i\ker b_i$ by~\eqref{eq:sibling-kernel}. Since $q_i-1\geq1$ and $h_i\geq1$, the complement is nonzero. Thus~\eqref{eq:cone-decomposition} extends each specified frame to a simplex of $X_r(R)$.
\end{proof}

\begin{corollary}\label{cor:frame-homology}
For $r\geq3$, one has $\widetilde H_d(X_r(R),\Z)=0$ for $0\leq d\leq r-3$.
\end{corollary}

\begin{proof}
Let $z$ be a finite ordered integral $d$-cycle, using the reduced chain complex when $d=0$. Its support, together with all faces of its simplices, is finite, and every frame in this family has at most $d+1\leq r-2$ vectors. Choose a common extension $v$ by Proposition~\ref{prop:frame-cone}. Prepending $v$ defines a chain operator $h_v$ on this finite family, and the ordered boundary satisfies $\partial h_v+h_v\partial=\id$. At the augmentation, set $h_v(1)=[v]$; the signs are the usual integral alternating signs. Since $\partial z=0$, one has $\partial h_vz=z$. The same augmented calculation for reduced zero-cycles proves connectedness.
\end{proof}

\begin{corollary}\label{cor:frame-simple-connectivity}
The realization $|X_4(R)|$ is connected and simply connected.
\end{corollary}

\begin{proof}
By Proposition~\ref{prop:frame-cone}, any two vertices admit a common apex and are therefore joined by a path, so $|X_4(R)|$ is connected. Let a finite edge loop be given. Choose one common apex $v$ for all its vertices and ordered edges, with complements for the resulting frames. For each ordered edge $(u,w)$ traversed by the loop, use the triangle $(v,u,w)$. Adjacent triangles share their radial edges and therefore form a continuous triangular fan whose boundary is the given loop. A reverse traversal uses the same ordered edge with the opposite orientation; it does not replace that edge by the distinct cell $(w,u)$. Repeated vertices cause no difficulty for the resulting map of a disk. By cellular approximation, every element of the fundamental group is represented by a finite edge loop, and the fan contracts it. Thus $|X_4(R)|$ is simply connected.
\end{proof}

\section{An acyclicity criterion and its application}
\label{sec:acyclicity-criterion}

Throughout this section, $A$ is a nonzero unital $\F_2$-algebra, $\Gamma_r=\GL_r(A)$, and $X_r(A)$ is the ordered frame semisimplicial set of Definition~\ref{def:ordered-frames}.

\begin{theorem}\label{thm:acyclicity-criterion}
Suppose that $A^2\cong A$ as right $A$-modules. Assume also that:
\begin{enumerate}
\item for every $n>0$, the standard inclusion $\Gamma_{n+2}\to\Gamma_{n+3}$, $g\mapsto\operatorname{diag}(g,1)$, induces zero on $H_n(-,\Z)$;
\item $\widetilde H_i(X_r(A),\Z)=0$ for $0\leq i\leq r-3$ and $r\geq4$.
\end{enumerate}
Then $H_n(A^\times,\Z)=0$ for every $n>0$.
\end{theorem}

To control integral homology and the coefficient modules below, we transfer Quillen's scalar weight calculation through reduction modulo two; see \cite[\S11, Lemmas~15--16]{Quillen72}.

\begin{lemma}\label{lem:scalar-homology}
Let $m\geq1$, let $F=\F_{2^m}$ and $C=F^\times$, and regard an $F$-vector space $V$ as a discrete additive group. For $1\leq j<m$, the group $H_j(V,\Z)$ is annihilated by $2$, and the norm element $N_C=\sum_{c\in C}[c]\in\Z[C]$ acts as zero on it. If a group $H$ acts $F$-linearly on $V$, then $N_C$ acts as zero on $H_i(H,H_j(V,\Z))$ for every $i\geq0$. These homology groups and all their $C$-stable subquotients have zero $C$-coinvariants. The same conclusions hold for inverse scalar multiplication.
\end{lemma}

\begin{proof}
Suppose first that $V$ is finite dimensional over $F$. Its additive group is a finite product of copies of $C_2$. The integral cellular chain complex of $BC_2=\mathbb{RP}^{\infty}$ has one copy of $\Z$ in each nonnegative degree, with differential $2$ in positive even degrees and zero in odd degrees. It is the direct sum of $\Z[0]$ and two-term complexes $\Z\xrightarrow{2}\Z$ in degrees $(2a,2a-1)$, $a\geq1$. Multiplication by $2$ on each two-term complex is null-homotopic: take the identity map from degree $2a-1$ to degree $2a$ as the homotopy.

The product chain complex is the tensor product of these complexes. Every summand other than $\Z[0]$ contains a two-term factor. Applying the preceding homotopy to one such factor, with its tensor sign, makes multiplication by $2$ null-homotopic on that summand; the cross terms cancel. Consequently $2H_j(V,\Z)=0$ for $j>0$. In the long exact sequence associated with $\Z\xrightarrow{2}\Z\to\F_2$, the kernel of reduction modulo two is the image of multiplication by $2$. Thus reduction induces a natural injection
$$
H_j(V,\Z)\longrightarrow H_j(V,\F_2)\qquad\text{for }j>0, \eqnum[eq:integral-mod-two]
$$
which commutes with every automorphism of $V$.

The mod-two cohomology of $V$ is the polynomial algebra on $H^1(V,\F_2)=\Hom_{\F_2}(V,\F_2)$. Indeed, $H^*(C_2,\F_2)=\F_2[t]$ with $|t|=1$, and this description extends to a finite product by the field K\"unneth formula. Since the algebra is generated by degree one, the description is natural under all linear automorphisms of $V$.

After extending scalars to a splitting field, the eigencharacters of scalar multiplication on $V$ are $\lambda\mapsto\lambda^{2^a}$, $0\leq a<m$, each with multiplicity $\dim_F V$. They may also be read from the isomorphism
$$
F\otimes_{\F_2}F\longrightarrow\prod_{a=0}^{m-1}F,\qquad x\otimes y\longmapsto(x^{2^a}y)_a.
$$
With the contragredient convention the cohomology generators have weights $-2^a$. Dualizing each finite cohomological degree shows that the weights on $H_j(V,\F_2)$ are sums of exactly $j$ members of $\{1,2,4,\ldots,2^{m-1}\}$, with repetitions, modulo $d=2^m-1$.

Quillen's counting lemma \cite[Lemma~16]{Quillen72}, specialized to $p=2$, shows that no such sum is zero modulo $d$ when $0<j<m$. Its cyclic carrying argument is as follows. Place one counter for each summand in the corresponding cyclic position of $\Z/m\Z$. Replace two counters in one position by one counter in the next. Each replacement preserves the residue modulo $d$, including the replacement in the last position, and strictly decreases the positive number of counters. The process terminates at a nonempty set of distinct positions with fewer than $m$ counters. Its ordinary sum lies between $1$ and $d-1$, so its residue is nonzero.

The order $d$ is odd, and the norm vanishes on every nontrivial character in characteristic two. By the weight calculation, $N_C$ acts as zero on $H_j(V,\F_2)$ for $1\leq j<m$, and hence on $H_j(V,\Z)$ by~\eqref{eq:integral-mod-two}. Inverse scalar multiplication replaces every weight by its negative and leaves the conclusion unchanged.

For arbitrary $V$, its finite-dimensional $F$-subspaces form a directed system of $C$-stable subgroups. Every finite bar chain is contained in one of them. Since filtered colimits of abelian groups are exact, the bar complex shows that group homology commutes with this union. The exponent and norm assertions therefore hold without a finiteness assumption on $V$.

If $H$ acts $F$-linearly, its action commutes with $C$. On the bar complex computing $H_i(H,H_j(V,\Z))$, the norm acts coefficientwise as zero, and hence induces zero on homology. These homology groups are annihilated by $2$, and both properties pass to $C$-stable subquotients. On coinvariants the norm acts as multiplication by the odd integer $d$, hence as the identity on a group annihilated by $2$. The coinvariants are therefore zero.
\end{proof}

If $C$ has odd order $d$ and $M$ is a $2$-primary $C$-module, multiplication by $d$ is invertible on $M$. The averaging operator $d^{-1}N_C$ identifies invariants and coinvariants and makes both functors exact on this category. Thus, if $M_C=0$, every $C$-stable subquotient of $M$ has zero coinvariants. For a module annihilated by $2^a$, the inverse of $d$ may be represented by an integer inverse modulo $2^a$. We will use this exactness for a kernel whose graded pieces are annihilated by $2$ but whose exponent is bounded by $2^n$.

Assume now that $A^2\cong A$. Then $A^m\cong A$ for every $m\geq1$, by induction using $A^{m+1}\cong A^m\oplus A$. Represent one such isomorphism by a row $L$ and its inverse by a column $L'$, so that $LL'=1$ and $L'L=I_m$. Then $T\mapsto LTL'$ and $a\mapsto L'aL$ are mutually inverse unital ring isomorphisms between $M_m(A)$ and $A$, and hence $\Gamma_m\cong A^\times$.

Composing the regular representation $\F_{2^m}\to M_m(\F_2)$ with the entrywise inclusion $M_m(\F_2)\hookrightarrow M_m(A)$ and the isomorphism $M_m(A)\to A$ embeds $\F_{2^m}$ into $A$. Indeed, the composite is unital and therefore injective because $A\neq0$.

For $k,q\geq1$, put
$$
J_{k,q}(A)=
 \left\{\begin{pmatrix}I_k&B\\0&H\end{pmatrix}:
 B\in M_{k,q}(A)\text{ and }H\in\Gamma_q\right\}. \eqnum[eq:frame-stabilizer]
$$
Write $U=M_{k,q}(A)_{\mathrm{add}}$. The extension $1\to U\to J_{k,q}(A)\xrightarrow{\pi}\Gamma_q\to1$ splits by $s(H)=\operatorname{diag}(I_k,H)$. If $u(B)=\left(\begin{smallmatrix}I_k&B\\0&I_q\end{smallmatrix}\right)$, then $s(H)u(B)s(H)^{-1}=u(BH^{-1})$. Thus the left $\Gamma_q$-action on $U$ is $B\mapsto BH^{-1}$. In particular,
$$
\begin{pmatrix}I_k&B\\0&H\end{pmatrix}
 =u(BH^{-1})s(H)\qquad\text{and}\qquad
 \begin{pmatrix}I_k&B\\0&H\end{pmatrix}^{-1}
 =\begin{pmatrix}I_k&-BH^{-1}\\0&H^{-1}\end{pmatrix}.
$$

Fix $n>0$, choose $m>n$, and embed $F=\F_{2^m}$ into $A$ as above. Let $C=F^\times$. Conjugation by $t_\lambda=\operatorname{diag}(\lambda I_k,I_q)$ fixes $s(\Gamma_q)$ pointwise and sends $B\in U$ to $\lambda B$.

The next calculation is a form of the scalar argument for affine groups; compare \cite[Theorem~1.9]{Suslin84}, \cite[Proposition~1.10, Theorem~1.11 and Remark~1.13]{NS90}, \cite[\S2]{Schlichting17}, and \cite[Proposition~2.1]{Knudson02}. The role of central scalars in the Nesterenko--Suslin argument is noted in \cite[Definition~2.1, footnote~1]{Schlichting17}. Here $F$ need not be central in $A$: left scalar multiplication commutes with the $\Gamma_q$-action, since $\lambda(BH^{-1})=(\lambda B)H^{-1}$.

\begin{lemma}\label{lem:stabilizer-scalars}
For the scalar action just defined, let
$$
Q_n=\ker\bigl(H_n(J_{k,q}(A),\Z)\xrightarrow{\pi_*}H_n(\Gamma_q,\Z)\bigr).
$$
Then $2^nQ_n=0$, $(Q_n)_C=0$, and the projection induces an isomorphism
$$
H_n(J_{k,q}(A),\Z)_C\cong H_n(\Gamma_q,\Z). \eqnum[eq:stabilizer-coinvariants]
$$
If the standard inclusion $\Gamma_q\to\Gamma_{k+q}$, $H\mapsto\operatorname{diag}(H,I_k)$, induces zero on $H_n(-,\Z)$, then the inclusion $J_{k,q}(A)\to\Gamma_{k+q}$ also induces zero on $H_n(-,\Z)$. All these assertions remain valid when the scalar actions are inverted.
\end{lemma}

\begin{proof}
Associated with the split extension is the $C$-equivariant Lyndon--Hochschild--Serre spectral sequence
$$
E^2_{p,j}=H_p(\Gamma_q,H_j(U,\Z)) \ \Longrightarrow\ H_{p+j}(J_{k,q}(A),\Z). \eqnum[eq:stabilizer-lhs]
$$
We use its usual homological indexing and projection edge map; see \cite[Theorem~6.8.2]{Weibel94}. The row $j=0$ survives unchanged to $E^\infty$. Indeed, the section and projection compare it with the spectral sequence of the extension of $\Gamma_q$ with trivial kernel and induce the identity on that row. There are no incoming differentials, and naturality with the section makes the outgoing differentials zero.

Consequently $Q_n=F_{n-1}H_n(J_{k,q}(A),\Z)$, with graded pieces $E^\infty_{p,n-p}$ for $0\leq p<n$. Each is a $C$-stable subquotient of $E^2_{p,n-p}$ in~\eqref{eq:stabilizer-lhs}, where $1\leq n-p\leq n<m$. By Lemma~\ref{lem:scalar-homology}, these pieces are annihilated by $2$ and have zero $C$-coinvariants. There are at most $n$ pieces, so $2^nQ_n=0$. Exactness of coinvariants on $2$-primary modules, applied successively to the filtration, implies $(Q_n)_C=0$.

The section is fixed by $C$ and splits $\pi_*$. Together with $(Q_n)_C=0$, this proves~\eqref{eq:stabilizer-coinvariants}. Let $i:J_{k,q}(A)\to\Gamma_{k+q}$ be the inclusion. If $\alpha_\lambda$ denotes the scalar automorphism of $J_{k,q}(A)$, then $i\alpha_\lambda=c_{t_\lambda}i$, where $c_{t_\lambda}$ is conjugation in $\Gamma_{k+q}$. Since inner automorphisms induce the identity on homology, $i_*$ is $C$-invariant and therefore factors through the coinvariant quotient $H_n(J_{k,q}(A),\Z)_C$. Under~\eqref{eq:stabilizer-coinvariants}, the resulting map is induced by $H\mapsto\operatorname{diag}(I_k,H)$. A block permutation conjugates this map to $H\mapsto\operatorname{diag}(H,I_k)$, the specified standard inclusion, whose map on $H_n$ is zero by hypothesis. Inverting the scalar actions permutes the summands of the norm and preserves the argument.
\end{proof}

\begin{proof}[Proof of Theorem~\ref{thm:acyclicity-criterion}]
Let $\mathcal B_r=E\Gamma_r\times_{\Gamma_r}|X_r(A)|$. For a free right $\Z\Gamma_r$-resolution $P_*$ of $\Z$, the double complex $P_t\otimes_{\Z\Gamma_r}C_p(X_r(A),\Z)$ computes $H_*(\mathcal B_r,\Z)$. The spectral sequence obtained by taking homology first in the frame direction has $E^2_{a,b}=H_a(\Gamma_r,H_b(X_r(A),\Z))$. By the second hypothesis, the projection $\mathcal B_r\to B\Gamma_r$ induces an isomorphism on $H_n$ whenever $r\geq n+3$.

For the other filtration, $\Gamma_r$ acts transitively on the ordered $(k-1)$-simplices: an isomorphism from $A^{r-k}$ onto the complement in~\eqref{eq:frame-decomposition} completes a given frame to an invertible matrix. The stabilizer of the standard $k$-frame is $J_{k,r-k}(A)$, and the stabilizer of a full frame is trivial; put $J_{r,0}(A)=1$. Since the frames are ordered, these stabilizers fix their simplices pointwise. By Shapiro's lemma, the spectral sequence is
$$
E^1_{p,t}=H_t(J_{p+1,r-p-1}(A),\Z) \ \Longrightarrow\ H_{p+t}(\mathcal B_r,\Z), \qquad\text{for }0\leq p\leq r-1; \eqnum[eq:frame-spectral-sequence]
$$
see \cite[\S6.1.15]{Weibel94} and compare \cite[Lemmas~2.3--2.4]{NS90}. All filtrations in a fixed total degree are finite.

The row $t=0$ in~\eqref{eq:frame-spectral-sequence} has one copy of $\Z$ in each column. The differential out of column $p$ is multiplication by $\sum_{i=0}^p(-1)^i$. Thus
$$
E^2_{p,0}=0\qquad\text{for }1\leq p\leq r-2. \eqnum[eq:frame-bottom-row]
$$
For positive even $p$ the outgoing differential is injective, and for odd $p$ the incoming differential from $p+1$ is the identity.

We induct on $n$. Assume $H_t(A^\times,\Z)=0$ for $0<t<n$; when $n=1$ this assumption is empty. Set $r=n+3$. The projection $\mathcal B_r\to B\Gamma_r$ then induces an isomorphism in degree $n$. Choose $m>n$ so that Lemma~\ref{lem:scalar-homology} applies in every positive degree at most $n$, and fix a unital copy of $F=\F_{2^m}$ in $A$. Put $C=F^\times$.

For $\lambda\in C$, define
$$
T_\lambda(v_1,\ldots,v_k)=(v_1\lambda,\ldots,v_k\lambda). \eqnum[eq:global-rescaling]
$$
Right multiplication by $\lambda$ need not be $A$-linear when $F$ is not central, but it preserves every frame: $v_i\lambda A=v_iA$, so the same complement may be used. The maps commute with deletion and satisfy $g(v_i\lambda)=(gv_i)\lambda$ for $g\in\Gamma_r$. Since $C$ is abelian, they define a $C$-action on the ordered semisimplicial set.

The identity on $E\Gamma_r$ together with~\eqref{eq:global-rescaling} defines an action on $\mathcal B_r$ covering the identity of $B\Gamma_r$. On the double complex, $1\otimes T_\lambda$ commutes with both differentials and preserves both filtrations. The Borel projection is therefore $C$-equivariant when the base has the trivial action. Since it induces an isomorphism on $H_n$, the $C$-action on $H_n(\mathcal B_r,\Z)$ is trivial.

To identify the induced action on~\eqref{eq:frame-spectral-sequence}, let $\sigma_k$ be the standard $k$-frame and put $t_{\lambda,k}=\operatorname{diag}(\lambda I_k,I_{r-k})$. Then $T_\lambda\sigma_k=t_{\lambda,k}\sigma_k$, so transport back to $\sigma_k$ by $t_{\lambda,k}^{-1}$ induces the stabilizer automorphism
$$
h\longmapsto t_{\lambda,k}^{-1}ht_{\lambda,k},\qquad\text{that is,}\qquad
 \begin{pmatrix}I_k&B\\0&H\end{pmatrix}
 \longmapsto\begin{pmatrix}I_k&\lambda^{-1}B\\0&H\end{pmatrix}.
$$
Equivalently, the orbit map $gJ\mapsto gt_{\lambda,k}J$ induces $x\mapsto xt_{\lambda,k}$ on $E\Gamma_r/J$, and $(xh)t_{\lambda,k}=(xt_{\lambda,k})(t_{\lambda,k}^{-1}ht_{\lambda,k})$ confirms the conjugation convention. Thus the action on the $E^1$-term is the inverse scalar action of Lemma~\ref{lem:stabilizer-scalars}.

The chosen deletion representatives commute with these stabilizer automorphisms only up to inner conjugation. If the rows of $B$ are $b_1,\ldots,b_k$, deleting row $i$ and placing that coordinate first among the new quotient coordinates defines
$$
f_i\begin{pmatrix}I_k&B\\0&H\end{pmatrix}
 =\begin{pmatrix}
 I_{k-1}&0&B_{\widehat i}\\
 0&1&b_i\\
 0&0&H
 \end{pmatrix}.
$$
The source action multiplies both $B_{\widehat i}$ and $b_i$ by $\lambda^{-1}$, whereas the target action multiplies only $B_{\widehat i}$. Their discrepancy is conjugation by $\operatorname{diag}(I_{k-1},\lambda^{-1},I_q)$ in the target stabilizer, where $q=r-k$. Inner automorphisms induce the identity on homology, and the global action on the double complex ensures compatibility with every page of the spectral sequence.

For $1\leq p<n$, put $t=n-p$ and $q=r-p-1$. Then $0<t<n$ and $q\geq3$. Since $\Gamma_q\cong A^\times$, it follows from the induction hypothesis that $H_t(\Gamma_q,\Z)=0$. By Lemma~\ref{lem:stabilizer-scalars}, with the inverse scalar action, $H_t(J_{p+1,q}(A),\Z)=Q_t$ is $2$-primary and has zero $C$-coinvariants. Each $E^\infty_{p,n-p}$ is a $C$-stable subquotient of $E^1_{p,n-p}$, so exactness of coinvariants on $2$-primary modules implies
$$
\bigl(E^\infty_{p,n-p}\bigr)_C=0\qquad\text{for }1\leq p<n. \eqnum[eq:lower-filtration-coinvariants]
$$
When $n=1$, there are no pieces in this range.

Since $r=n+3$, \eqref{eq:frame-bottom-row} implies $E^\infty_{n,0}=0$. For $p=0$, the stabilizer map and the Borel projection form
$$
H_n(J_{1,r-1}(A),\Z)\longrightarrow H_n(\mathcal B_r,\Z)\xrightarrow{\cong}H_n(\Gamma_r,\Z).
$$
The composite is induced by the stabilizer inclusion and is zero by hypothesis~(1) and Lemma~\ref{lem:stabilizer-scalars}, since $r-1=n+2$. The first map has image $F_0H_n(\mathcal B_r,\Z)$; hence this filtration subgroup is zero.

All graded pieces of the finite filtration of $H_n(\mathcal B_r,\Z)$ now have zero $C$-coinvariants. Applying the right-exact coinvariant functor successively to the short exact sequences of this filtration, we obtain $H_n(\mathcal B_r,\Z)_C=0$. Since the $C$-action on $H_n(\mathcal B_r,\Z)$ is trivial, $H_n(\mathcal B_r,\Z)=0$. The Borel comparison and the isomorphism $\Gamma_r\cong A^\times$ then imply $H_n(A^\times,\Z)=0$, completing the induction.
\end{proof}

\begin{theorem}\label{thm:acyclicity}
The unit group of $R=L_{\F_2}(1,2)$ is integrally acyclic. More generally, $H_n(\GL_r(R),\Z)=0$ for every $n>0$ and $r\geq1$.
\end{theorem}

\begin{proof}
Lemma~\ref{lem:leaf-coordinates}, Proposition~\ref{prop:zero-padding}, and Corollary~\ref{cor:frame-homology} verify the hypotheses of Theorem~\ref{thm:acyclicity-criterion}. It follows that $H_n(R^\times,\Z)=0$ for all $n>0$. The leaf ring isomorphisms~\eqref{eq:leaf-ring-isomorphisms} identify every $\GL_r(R)$ with $R^\times$, so the same conclusion holds in every rank.
\end{proof}

\section{The Steinberg comparison}
\label{sec:steinberg}

The acyclicity of the unit group will be used to verify the following criterion for the Steinberg map.

For a nonzero unital ring $B$, write $t_{ij}(a)=I+aE_{ij}$ and let $E_m(B)$ be the subgroup of $\GL_m(B)$ generated by these matrices. For $m\geq3$, the ordinary Steinberg group $\St_m(B)$ has generators $X_{ij}(a)$, where $i\ne j$ and $a\in B$, and relations
$$
\begin{aligned}
\eqrownum &\quad X_{ij}(a)X_{ij}(b)=X_{ij}(a+b),\\
\eqrownum &\quad [X_{ij}(a),X_{kl}(b)]=1 \quad \text{for }i\ne l\text{ and }j\ne k,\\
\eqrownum &\quad [X_{ij}(a),X_{jk}(b)]=X_{ik}(ab) \quad \text{for distinct }i,j,k.
\end{aligned}
$$
Here $[x,y]=xyx^{-1}y^{-1}$. Let $\phi_m:\St_m(B)\to\GL_m(B)$ send $X_{ij}(a)$ to $t_{ij}(a)$, and put $N_m(B)=\ker\phi_m$. The standard stabilization homomorphism $j_m:\St_m(B)\to\St_{m+1}(B)$ is defined by $j_m(X_{ij}(a))=X_{ij}(a)$.

For a ring $B$ of characteristic two, let $X_n(B)$ be the ordered frame semisimplicial set of Definition~\ref{def:ordered-frames}.


\begin{theorem}[Steinberg comparison criterion]
\label{thm:raw-criterion}
Let $B$ be a nonzero unital ring of characteristic two, and fix $n\geq4$. Suppose that
\begin{enumerate}
\item $\GL_{n-1}(B)=E_{n-1}(B)$ and $\GL_{n-2}(B)=E_{n-2}(B)$;
\item $j_{n-1}(N_{n-1}(B))=1$ in $\St_n(B)$;
\item $|X_n(B)|$ is simply connected.
\end{enumerate}
Then $\phi_n:\St_n(B)\to\GL_n(B)$ is an isomorphism.
\end{theorem}

\begin{proof}
Put $G=\GL_n(B)$, $S=\St_n(B)$, and let $b_1,\ldots,b_n$ be the standard columns.

The action of $G$ on ordered $k$-frames is transitive, since a basis of a permitted complement completes the frame to an automorphism of $B^n$. In particular, there is one orbit in each of dimensions zero, one, and two. The vertex stabilizer $J$ and the stabilizer $K$ of the ordered edge are
$$
J=\left\{j(b,H)=\begin{pmatrix}1&b\\0&H\end{pmatrix}\right\}\qquad\text{and}\qquad
 K=\left\{k(a,b,H)=
 \begin{pmatrix}1&0&a\\0&1&b\\0&0&H\end{pmatrix}\right\}.
$$
In $J$, the row $b$ has length $n-1$ and $H\in\GL_{n-1}(B)$. In $K$, the rows have length $n-2$ and $H\in\GL_{n-2}(B)$. The upper rows in these full stabilizers are unrestricted. Let $\tau=(12)$ and $h=(23)$ be permutation matrices, so that $h\in J$, and put $\eta(k)=\tau k\tau^{-1}$. This automorphism of $K$ interchanges its two upper rows.

We apply Brown's presentation for a group acting on a simply connected complex \cite[Theorem~1 and \S3]{Brown84}. For the ordered frame complex, the Borel construction $Y=EG\times_G|X_n(B)|$ has fundamental group
$$
\Pi=\langle J,T\mid TkT^{-1}=\eta(k)\text{ for }k\in K,\quad ThT=hTh\rangle. \eqnum[eq:borel-presentation]
$$
This notation includes every multiplication relation in $J$. To see that the remaining relations are complete, consider the cellular Borel construction through total dimension two. The vertex orbit contributes $BJ$. In the cylinder over $BK$, the product of the zero-cell of $BK$ with the edge has total dimension one and supplies $T$; the products of the one-cells of $BK$ with the edge impose the conjugation relations between the endpoint inclusions. The zero-cell in the classifying space of the triangle stabilizer contributes one further two-cell. All other cells have total dimension at least three.

To determine the triangle relation, identify $\pi_1(Y)$ with the group of pairs $(g,[p])$, where $p$ is a path from $b_1$ to $gb_1$ taken up to homotopy relative to its endpoints, with multiplication
$$
(g,[p])(g',[p'])=(gg',[p*(gp')]).
$$
Represent $j\in J$ by its constant path and $T$ by $(\tau,E)$, where $E$ is the ordered edge $(b_1,b_2)$. For $k\in K$, both $Tk$ and $\eta(k)T$ have projection $\tau k$ and path $E$. The word $ThT$ has projection $(13)$ and path $b_1\to b_2\to b_3$, whereas $hTh$ has the same projection and path $b_1\to b_3$. The boundary of the ordered triangle $(b_1,b_2,b_3)$ therefore imposes the second relation in \eqref{eq:borel-presentation}.

The cells $(u,v)$ and $(v,u)$ are distinct; traversing one in reverse does not replace it by the other. Since $h^2=1$, the braid relation implies $(hT)h(hT)^{-1}=T$ and hence $T^2=1$.

The homotopy sequence of $|X_n(B)|\to Y\to BG$, together with the third hypothesis, shows that the projection $p:\Pi\to G$, defined by $p(j)=j$ and $p(T)=\tau$, is an isomorphism.

We now lift this presentation to $S$. Permuting indices is an automorphism of the Steinberg presentation, so the second hypothesis also applies to stabilization on coordinates $2,\ldots,n$. For $H\in\GL_{n-1}(B)=E_{n-1}(B)$, choose an elementary word representing $H$ and evaluate it on these coordinates. If two words represent the same $H$, their quotient is a word in $N_{n-1}(B)$; the second hypothesis makes its value in $S$ trivial. Hence the evaluation is independent of the chosen word, and concatenation defines a homomorphism
$$
\ell:\GL_{n-1}(B)\longrightarrow S
$$
covering $H\mapsto\operatorname{diag}(1,H)$.

For a row $b=(b_2,\ldots,b_n)$ set $x_1(b)=\prod_{j=2}^nX_{1j}(b_j)$, in increasing order. The factors commute. According to the Steinberg relations,
$$
\ell(H)x_1(b)\ell(H)^{-1}=x_1(bH^{-1}). \eqnum[eq:row-action]
$$
Indeed, for $H=t_{ij}(c)$ with $i,j\geq2$, conjugation changes only the $j$-th entry, replacing it by $b_j-b_i c$. This is right multiplication by $I-cE_{ij}$; the third Steinberg relation produces the new term, while all other factors commute. Elementary generation extends the identity to every $H$.

Define $\sigma(j(b,H))=\ell(H)x_1(b)$. By \eqref{eq:row-action},
$$
\sigma(j(b,H))\sigma(j(b',H')) =\ell(HH')x_1(bH'+b') =\sigma(j(bH'+b',HH')),
$$
which is precisely the matrix multiplication law in $J$. Thus $\sigma:J\to S$ is a homomorphism lifting the inclusion $J\to G$.

For $H\in\GL_{n-2}(B)$, set $\ell_0(H)=\ell(\operatorname{diag}(1,H))$. Since $\GL_{n-2}(B)=E_{n-2}(B)$, an elementary word representing $H$ uses only coordinates $3,\ldots,n$. Its evaluation is $\ell_0(H)$ by the definition of $\ell$, so this lift has support on those coordinates. For $n=4$, the word is evaluated through $\St_3(B)$ on coordinates $2,3,4$, without using a rank-two Steinberg group. For $k(a,b,H)\in K$, one consequently has
$$
\sigma(k(a,b,H))=\ell_0(H)x_2(b)x_1(a), \eqnum[eq:edge-section]
$$
where both row products range over columns $3,\ldots,n$.

Put $w_{ij}=X_{ij}(1)X_{ji}(1)X_{ij}(1)$. The Steinberg relations and characteristic two give $w_{ij}^2=1$. We next verify, inside the Steinberg presentation, that
$$
w_{12}X_{ij}(c)w_{12}^{-1} =X_{\tau(i),\tau(j)}(c). \eqnum[eq:weyl-action]
$$
Write $u=X_{12}(1)$ and $v=X_{21}(1)$, so that $w_{12}=uvu$. For $j\geq3$, the Steinberg relations imply
$$
\begin{aligned}
uX_{1j}(c)u^{-1}&=X_{1j}(c),\\
vX_{1j}(c)v^{-1}&=X_{2j}(c)X_{1j}(c),\\
u\bigl(X_{2j}(c)X_{1j}(c)\bigr)u^{-1}&=X_{2j}(c),
\end{aligned}
\qquad\text{and}\qquad
\begin{aligned}
uX_{j1}(c)u^{-1}&=X_{j2}(c)X_{j1}(c),\\
v\bigl(X_{j2}(c)X_{j1}(c)\bigr)v^{-1}&=X_{j2}(c),\\
uX_{j2}(c)u^{-1}&=X_{j2}(c).
\end{aligned}
$$
The cancellations in the last terms use characteristic two. Since $w_{12}^2=1$, the same formulas establish the identities with $1$ and $2$ interchanged. Generators $X_{ij}(c)$ with $i,j\notin\{1,2\}$ commute with both $u$ and $v$. For the opposite roots, choose $j\geq3$ and use
$$
X_{12}(c)=[X_{1j}(c),X_{j2}(1)].
$$
The preceding formulas send the right-hand side to $[X_{2j}(c),X_{j1}(1)]=X_{21}(c)$, and the reverse case follows again from $w_{12}^2=1$. Thus \eqref{eq:weyl-action} holds for every defining generator.

The element $w_{12}$ commutes with $\ell_0(H)$ and interchanges the two row products in \eqref{eq:edge-section}; those products commute with each other. Hence
$$
w_{12}\sigma(k)w_{12}^{-1}=\sigma(\eta(k))
\qquad\text{for every }k\in K.
$$
Moreover, $\sigma(h)=w_{23}$, since the chosen elementary word on the last coordinates represents $(23)$. Applying \eqref{eq:weyl-action} and its version with indices permuted to the defining words for $w_{23}$ and $w_{12}$, we obtain
$$
w_{12}w_{23}w_{12}=w_{13}=w_{23}w_{12}w_{23},
$$
so the triangle relation also lifts to $S$.

The presentation \eqref{eq:borel-presentation} now defines a homomorphism $\psi:\Pi\to S$ with $\psi|_J=\sigma$, $\psi(T)=w_{12}$, and $\phi_n\psi=p$. It remains to prove that $\psi$ is surjective. Its image contains $X_{ij}(c)$ for $i,j\geq2$ through $\ell$, and contains $X_{1j}(c)$ through the upper-row subgroup of $\sigma(J)$. By \eqref{eq:weyl-action}, the conjugates of $X_{i2}(c)$ for $i\geq3$ and of $X_{12}(c)$ by $w_{12}$ are $X_{i1}(c)$ and $X_{21}(c)$, respectively. Thus $\psi$ is surjective.

Since $p$ is an isomorphism, $s=\psi p^{-1}:G\to S$ is a surjective section of $\phi_n$. It follows that $s\phi_n=1_S$, so $s$ is the inverse of $\phi_n$.
\end{proof}

\begin{thebibliography}{99}
\bibitem[AAS17]{AAS17} G.~Abrams, P.~Ara and M.~Siles Molina, \emph{Leavitt Path Algebras}, Lecture Notes in Mathematics, vol.~2191, Springer, London, 2017. \url{https://doi.org/10.1007/978-1-4471-7344-1}.
\bibitem[ABC09]{ABC09} P.~Ara, M.~Brustenga and G.~Corti\~nas, \emph{$K$-theory of Leavitt path algebras}, M\"unster J. Math. \textbf{2} (2009), 5--34. \url{https://arxiv.org/abs/0903.0056v2}.
\bibitem[Brown84]{Brown84} K.~S.~Brown, \emph{Presentations for groups acting on simply-connected complexes}, J. Pure Appl. Algebra \textbf{32} (1984), 1--10. \url{https://doi.org/10.1016/0022-4049(84)90009-4}.
\bibitem[Knudson02]{Knudson02} K.~P.~Knudson, \emph{Unstable homotopy invariance for finite fields}, Fund. Math. \textbf{175} (2002), no.~2, 155--162. \url{https://doi.org/10.4064/fm175-2-5}.
\bibitem[NS90]{NS90} Yu.~P.~Nesterenko and A.~A.~Suslin, \emph{Homology of the full linear group over a local ring, and Milnor's $K$-theory}, Math. USSR-Izv. \textbf{34} (1990), no.~1, 121--145. \url{https://doi.org/10.1070/IM1990v034n01ABEH000610}.
\bibitem[Quillen72]{Quillen72} D.~Quillen, \emph{On the cohomology and $K$-theory of the general linear groups over a finite field}, Ann. of Math. (2) \textbf{96} (1972), no.~3, 552--586. \url{https://doi.org/10.2307/1970825}.
\bibitem[Schlichting17]{Schlichting17} M.~Schlichting, \emph{Euler class groups and the homology of elementary and special linear groups}, Adv. Math. \textbf{320} (2017), 1--81. \url{https://doi.org/10.1016/j.aim.2017.08.034}.
\bibitem[Suslin84]{Suslin84} A.~A.~Suslin, \emph{Homology of $\mathrm{GL}_n$, characteristic classes and Milnor $K$-theory}, Trudy Mat. Inst. Steklov. \textbf{165} (1984), 188--204 (Russian); English translation, Proc. Steklov Inst. Math. \textbf{165} (1985), 207--226. \url{https://www.mathnet.ru/eng/tm2280}.
\bibitem[Wagoner72]{Wagoner72} J.~B.~Wagoner, \emph{Delooping classifying spaces in algebraic $K$-theory}, Topology \textbf{11} (1972), 349--370. \url{https://doi.org/10.1016/0040-9383(72)90031-6}.
\bibitem[Weibel13]{Weibel13} C.~A.~Weibel, \emph{The $K$-book: An Introduction to Algebraic $K$-theory}, Graduate Studies in Mathematics, vol.~145, American Mathematical Society, Providence, RI, 2013. \url{https://doi.org/10.1090/gsm/145}.
\bibitem[Weibel94]{Weibel94} C.~A.~Weibel, \emph{An Introduction to Homological Algebra}, Cambridge Studies in Advanced Mathematics, vol.~38, Cambridge University Press, Cambridge, 1994. \url{https://doi.org/10.1017/CBO9781139644136}.
\bibitem[vdK80]{vdK80} W.~van der Kallen, \emph{Homology stability for linear groups}, Invent. Math. \textbf{60} (1980), no.~3, 269--295. \url{https://doi.org/10.1007/BF01390018}.
\end{thebibliography}
\end{document}
